\section{Related work}\label{sec:related}

The problem sits where several literatures meet: the statistical comparison of classifiers, the label-efficient evaluation of machine learning models, the design of studies that compare diagnostic tests, the estimation of fairness metrics, prediction-assisted inference and classical survey sampling. We review each in turn, stating what the work does and how our setting differs, and close with a short positioning statement.

\subsection{Statistical evaluation and comparison of classifiers}

The comparison of two classifiers on a common test set is usually treated as a paired problem. McNemar's test \citep{mcnemar1947} uses only the discordant pairs of a two-by-two table of correct and incorrect classifications, and \citet{dietterich1998} studied its behaviour, together with resampled and cross-validated $t$ tests, as a tool for comparing learning algorithms. For areas under correlated ROC curves the nonparametric approach of \citet{delong1988} is standard. \citet{demsar2006} recommended rank-based tests for comparing several classifiers over several data sets. These methods share two assumptions that do not hold here. The test set is taken as fully labelled, so the question of which instances to label does not arise, and the metric is overall accuracy or a ranking measure, not a difference between models in group-specific rates. The use of discordant pairs in McNemar's test is, however, the simplest instance of the first structural fact on which our design rests: for a paired comparison of accuracy, concordant instances carry no information about the difference.

Group-fairness criteria that depend on the outcome, such as equal opportunity and equalised odds \citep{hardt2016}, add group-specific denominators to every comparison, and it is known that several such criteria cannot in general be satisfied at once when base rates differ between groups \citep{chouldechova2017}. We do not take a position on which criterion is appropriate; we take the true and false positive rate gaps and error as given and ask how to compare several models on them with few labels.

\subsection{Active testing and active evaluation of one model}

A large literature reduces the labelling cost of estimating the performance of a single trained model. Active testing \citep{kossen2021active} selects test points sequentially with acquisition functions tailored to evaluation and corrects the resulting bias with an importance-weighted estimator; active surrogate estimators \citep{kossen2022ase} instead interpolate the loss with an actively trained surrogate, and the approach has been scaled to large language models with cheap in-context surrogates \citep{berrada2025}. Prediction-powered active testing \citep{ashouritaklimi2026} combines the importance-weighted estimator with a prediction-powered control variate, derives variance-minimising acquisition rules and proves asymptotic normality. Bayesian assessment of black-box classifiers selects instances by Thompson-type rules \citep{ji2021active}. Importance sampling designs for F-measures and other ratio metrics were derived by \citet{sawade2010fmeasure} and, for rare categories and extreme class imbalance, by \citet{poms2021} and \citet{marchant2017}; \citet{yilmaz2021} compared importance and Poisson sampling designs for accuracy and F1 with confidence intervals.

Closer to our design are the stratified approaches. \citet{bennett2010} stratified on classifier scores to estimate the precision of deployed classifiers with an online approximation to optimal allocation, \citet{katariya2012} adaptively refined a stratification for accuracy estimation, and \citet{kumar2018} studied stratification on classifier output for risk estimation and reported large variance reductions against simple random sampling. More recently, approximate Neyman allocation from surrogate signals has been used to stratify evaluation pools for generative models \citep{liu2026neyman}, and stratified importance sampling with one-shot batch acquisition has been proposed for monitoring deployed classifiers \citep{marsigli2026}. All of this work estimates the performance of one model, usually a mean or an overall ratio such as precision, and builds strata from that model's scores. Our strata come from the agreement pattern of several models crossed with the protected group, our targets are between-model contrasts in group-specific rates, and our allocation must serve many such contrasts at once. Our design is also non-sequential: the whole budget is allocated before any label is seen, which suits labelling campaigns run as a single batch.

\subsection{Label-efficient model selection and comparison}

When several models are available, labels can be chosen to discriminate between them rather than to estimate each one well. \citet{sawade2012active} is the closest precedent for our use of disagreement. They compare the risks of two fixed models under a labelling budget, derive the instrumental distribution that maximises the power of a test of the risk difference, observe that it favours instances on which the models disagree, include a baseline that samples only where the models' predictions differ, and extend the approach to more than two models through a heuristic mixture of pairwise distributions, using the models' own predictive distributions as proxies for the unknown labels. \citet{hara2024} develop an adaptive labelling strategy that minimises the variance of estimated differences in test loss for selecting the best of several candidates. Consensus-driven active model selection \citep{kay2025coda} uses agreement and disagreement among candidate models in a probabilistic model to decide which points to label and updates a posterior over the identity of the best model. Discordant pair analysis \citep{musgrove2023} evaluates a new classifier against a baseline of known performance by adjudicating only the cases on which they differ, and \citet{kuboyama2026} shows that labelling every row on which at least two candidates disagree settles the accuracy ranking of several models exactly on a finite pool. Other work selects the best of several pretrained models from a stream or a pool with few labels \citep{karimi2021,okanovic2025}, chooses instances for annotation from clusters of output differences between two text generators \citep{ashurytahan2024}, or studies whether a fixed budget should be spent on more items or on repeated labels when identifying the more accurate of two classifiers \citep{dorner2024}. The only work we found that selects among several models on both accuracy and subgroup fairness under label costs is the Bayesian bandit formulation of \citet{kher2026}, in which an online policy decides which agent to query and the target is regret.

This literature establishes that disagreement is where comparative information lies, and we cite it as such. It differs from our setting in three respects. Its target is a risk or an accuracy ranking rather than a contrast of ratios with group-specific denominators; most of it is adaptive or Bayesian and aims at identifying a best model, not at estimating contrasts with design-based confidence intervals; and the designs that label disagreements only, such as \citet{musgrove2023} and \citet{kuboyama2026}, rely either on known baseline performance or on exhaustive labelling of the disagreement set. For error rates the denominator is the pool size and disagreements suffice, but for true and false positive rates the group prevalences $P_g$ and $Q_g$ depend on the agreement strata too.

\subsection{Verification designs for comparing diagnostic tests}

Biostatistics has long faced the problem of comparing two screening tests when the reference standard is too costly or invasive to apply to everyone. \citet{schatzkin1987} proposed applying the reference procedure only to people positive on at least one of the two tests and comparing the tests through ratios of sensitivities and of false positive rates; \citet{chock1997} studied the related design in which individuals negative on both tests are not verified; and \citet{pepe2001comparing} showed that, under verification of screen positives only, absolute sensitivities and specificities are not identified but their ratios are, and gave a regression framework for relative accuracy. Correcting for verification that depends on the test result is the subject of a classic literature on verification bias \citep{begg1983}, and the practice of resolving only discordant results with the reference standard, so-called discrepant analysis, has been criticised because it assumes that concordant results are correct \citep{hadgu1996}. Two-phase designs with a cheap first-phase classification and an optimally allocated second-phase verification were developed for the sensitivity, specificity and predictive value of one test by \citet{mcnamee2002}, within the broader two-phase framework of \citet{neyman1938} and its epidemiological descendants \citep{breslow1988,reilly1995}.

Our true and false positive rate contrasts between two classifiers are this design problem transplanted to machine learning, and the debt should be acknowledged. The differences are in scope and in what is identified. The diagnostic literature compares two tests, not several models; it does not have protected groups or contrasts of group gaps; its strata are the test results themselves, not results combined with predicted probabilities used to set the allocation; and several of its designs leave the doubly negative cell unverified on purpose, so that only relative measures are identified. Our design samples every stratum with known positive probability, so that differences in rates, not only their ratios, are identified, at the price of spending some labels where all models agree. How many such labels are worth spending is one of the allocation questions the design answers.

\subsection{Fairness estimation and auditing with few labels}

Several lines of work estimate fairness quantities with limited data. One line audits a model by querying it rather than by acquiring ground truth: active fairness auditing estimates demographic parity with few queries to the audited model \citep{yan2022active}, and query-efficient auditing of language models gives intervals for metrics such as group differences in AUC with active selection of queries, reporting large savings over stratified sampling \citep{hartmann2026}. Demographic parity needs no outcome labels, so the query-based setting is different from ours, although the second of these papers addresses an outcome-dependent metric for one model. A second line provides valid inference: \citet{chugg2023} give sequential, anytime-valid tests of fairness that allow data to be collected by changing and group-dependent policies, \citet{cherian2023} construct simultaneous bounds on performance disparities over many subgroups of one fully labelled model, \citet{ferrara2025} adapt hypothesis tests of fairness to group size, and \citet{lum2022} show that common disparity summaries are biased estimators when groups are small and propose corrections. A third line uses unlabelled data to improve estimates of group-specific rates. \citet{ji2020fairness} calibrate scores per group in a hierarchical Bayesian model to infer group differences in accuracy, true positive rate and false positive rate of one model, and remark that stratified sampling by group could help but is not always possible. \citet{gao2025} impute missing outcomes and augment for robustness to estimate group-specific true and false positive rates and calibration with semi-supervised efficiency gains, \citet{gronsbell2022} evaluate one prediction rule semi-supervisedly when the labels come from a stratified sample, \citet{miller2021} estimate subpopulation performance with model-based metrics, and \citet{das2025} audit fairness under partial feedback with a cost model for acquiring extra labels. Closest in setting is SSME \citep{shanmugam2025}, which fits a semi-supervised mixture model to the labels and the scores of several classifiers and estimates any metric for each classifier, including within demographic subgroups.

These methods either take the labelled set as given, so there is no design to optimise, or concern one model, or both. Our contribution is complementary: a sampling design and design-based estimators for between-model fairness contrasts. Model-based and semi-supervised estimators could in principle be applied to the labels our design collects, with design weights.

\subsection{Prediction-powered and model-assisted inference}

Prediction-powered inference \citep{angelopoulos2023ppi} combines a small labelled sample with machine learning predictions on a large unlabelled sample and gives confidence intervals that are valid whatever the quality of the predictions; PPI++ \citep{angelopoulos2026ppipp} tunes the weight on the predictions so that the interval is never wider than the classical one asymptotically. Stratified prediction-powered inference \citep{fisch2024stratppi} applies the idea within strata and allocates the label budget across strata by a Neyman-type rule. Active statistical inference \citep{zrnic2024active} samples labels with probabilities driven by model uncertainty and combines the resulting design-unbiased estimator with prediction-powered corrections; robust sampling \citep{li2025robust} interpolates between uniform and active sampling so that a poor proxy does not make matters worse than uniform sampling, and fixed-size variants replace Poisson sampling by conditional Poisson sampling \citep{skalnes2026}. Prediction-powered ranking \citep{chatzi2024} gives simultaneous rank-sets for several language models from pairwise comparisons. In these methods predictions enter either as a control variate in the estimator or as a guide to inclusion probabilities, and the target is typically a single mean or a vector of means for one model. We use predictions in the second way only, to set the allocation, and keep the estimator purely design-based; the target is a vector of nonlinear contrasts between models, and the strata are defined by the models' agreement pattern and the group.

The same idea is much older in survey sampling, where auxiliary information is used to build the design and to sharpen the estimator. Model-assisted estimation, in which a working model improves efficiency while validity rests on the design \citep{sarndal1992,breidt2017}, and calibration estimators \citep{deville1992} are the classical forms, and prediction-powered inference can be read as a member of this family. A model-assisted version of our estimator, which would use the models' probabilities also in the estimation step, is a natural extension that we do not pursue here.

\subsection{Survey-sampling foundations}

The design is assembled from standard survey-sampling components. Stratified sampling with allocation proportional to stratum size times stratum standard deviation, which minimises the variance of an estimated total for a fixed sample size, is due to \citet{neyman1934}. When several quantities are estimated from one sample, a single allocation must compromise between them; optimum allocation in multivariate surveys has a long history, including analytical solutions \citep{kokan1967}, and the weighted generalised Neyman rule we use minimises a weighted sum of variances, the simplest such compromise. Unbiased estimation of totals under unequal inclusion probabilities goes back to \citet{horvitz1952}. Group rates are ratios of totals, and the standard treatment of ratio estimators and of smooth functions of totals by Taylor linearisation is given by \citet{woodruff1971}, \citet{binder1983}, \citet{sarndal1992} and \citet{wolter2007}, with a practical account in \citet{lumley2010}. Two-phase sampling, in which a cheap first-phase measurement guides the selection of a second-phase sample, was introduced by \citet{neyman1938}; our setting is a two-phase design in which the first phase, the models' predictions on the whole pool, costs nothing. What these foundations do not supply is the choice of strata and of the variance targets for comparing classifiers on fairness, which is the part we work out.

\subsection{Positioning}

Disagreement-based sampling for comparing models, discordant-result designs for comparing diagnostic tests, proxy-based optimal stratification for evaluating one classifier, and label-efficient fairness estimation for one model are all established, and we build on each. To our knowledge they have not been combined into a fixed-budget design for estimating fairness contrasts between several classifiers. Such a design must treat group-prevalence denominators that depend on all strata, allocate labels for many contrasts at once, and provide design-based confidence intervals. That combination, its design-based theory and an empirical assessment of when it pays off are the subject of this paper.
